%----------------------------------------------------------------------------%
% modelo_relatorio.tex
% Copyright 2017 Elias Alves @unixelias
%
% This work may be distributed and/or modified under the
% conditions of the LaTeX Project Public License, either version 1.3
% of this license or (at your option) any later version.
% The latest version of this license is in
%   http://www.latex-project.org/lppl.txt
% and version 1.3 or later is part of all distributions of LaTeX
% version 2005/12/01 or later.
%
% This work has the LPPL maintenance status `maintained'.
% 
% The Current Maintainer of this work is Elias Alves @unixelias
%
% This work consists of the files modelo_relatorio.dtx
%
%%%%%%%%%%%%%%%%%%%%%%%%%%%%%%%%%%%%%%%%%%
%% Template para Apresentações tipo PPT
%% LaTeX Template
%% Version 1.0 (20/12/2017)
%%%%%%%%%%%%%%%%%%%%%%%%%%%%%%%%%%%%%%%%%%

\documentclass[12pt]{article} 
\usepackage[a4paper,left=30mm,right=25mm,top=30mm,bottom=25mm,headsep=15.5mm]{geometry} 

\usepackage[brazilian]{babel}                                
\usepackage[T1]{fontenc}
\usepackage[utf8]{inputenc}

\usepackage{enumitem}
\usepackage{ragged2e}

\usepackage{ae} 
\usepackage{amsfonts}
\usepackage{amssymb}
\usepackage{multicol}
\usepackage{pgf,pgfarrows,pgfnodes,pgfautomata,pgfheaps,pgfshade}
\usepackage{verbatim}
\usepackage{graphicx} 	% para inclusão de imagens
\usepackage{url}
\usepackage{minted} 	% inclusão de exemplos de código
\definecolor{fundo-cod}{rgb}{0.95,0.95,0.95}
\definecolor{exemplo}{HTML}{EAF3FA}
% ---
% para tratar urls longas
\usepackage{url}				% para tratar URLs longas
\makeatletter
\g@addto@macro{\UrlBreaks}{\UrlOrds}
\makeatother

% ---
% Pacotes para inserir licença do documento
% ---
\usepackage[
	type={CC},
	modifier={by-sa},
	version={4.0},
	]{doclicense}


\usepackage{indentfirst,graphicx} 
\usepackage[absolute]{textpos}


%----------------- ENTRADA DE DADOS -----------------------------
%================================================================
%---- Altere esses dados para preencher capa e folha de rosto ---
%================================================================
\def\titulotrabalho{Título do Trabalho}
\def\autores{Autor Um\\Autor Dois\\Autor Três}
\def\docente{Elias Alves}
\def\disciplina{Bancos de Dados II}
\def\tema{Tema para o seu trabalho}
\def\tipotrabalho{Trabalho prático?}
\def\dataentrega{02/01/2017}
\def\universidade{Universidade Federal dos Vales do Jequitinhonha e Mucuri}
\def\faculdade{Faculdade de Ciências Exatas e Tecnológicas - FACET}
\def\departamento{Departamento de Computação - Decom}
\def\curso{Sistemas de Informação}
\def\cidade{Diamantina}
\def\estado{MG}

%--------------------------- CAPA -------------------------------
%================================================================
%------------------------- Cabeçalho ----------------------------
%================================================================
\begin{document}
\pagestyle{myheadings}
\newpage
\pagenumbering{arabic} 
\setcounter{page}{1}
\renewcommand{\thepage}{\arabic{page}}
\thispagestyle{empty}

\begin{center}
    \begin{textblock*}{\linewidth}(2cm,1cm)%{width}[xt,yt](X,Y)
	    \raggedright
	    \includegraphics[width=0.17\linewidth]{../img/main/armas_brasil.eps}
    \end{textblock*} ~\\[1em]
	\begin{textblock*}{\linewidth}(4cm,1cm)%{width}[xt,yt](X,Y)
		\raggedleft
		\includegraphics[width=0.17\linewidth]{../img/main/logo_ufvjm.eps}
	\end{textblock*} ~\\[1em]
	\begin{textblock*}{0.73\linewidth}(5.1cm,1.4cm)%{width}[xt,yt](X,Y)
		\centering
			\bfseries\small\MakeUppercase{\universidade} \par
			\bfseries\scriptsize\MakeUppercase{\faculdade} \par
			\bfseries\scriptsize\MakeUppercase{\departamento} \par
			\bfseries\scriptsize\MakeUppercase{\curso} \par
	\end{textblock*} ~\\[1.5em]
\end{center}
%% Fim do cabeçalho
%---------------------------
\par %
\vfill
\vfill
\begin{center}
	\MakeUppercase\disciplina
\end{center}
\vfill
\vfill
\vfill
\begin{flushright}
	\autores
\end{flushright}
\vfill
\begin{center}
	\cidade{ }- Dezembro de 2017 
\end{center}
\clearpage
%================================================================
%------------------------- Folha de rosto -----------------------
%================================================================
\thispagestyle{empty}
\begin{center}
	\MakeUppercase{\universidade}
\end{center}
\vfill
\begin{center}
	\small\bfseries\MakeUppercase{\tema}
\end{center}
\vfill
\begin{flushright}
	\autores
\end{flushright}
\par
\begin{flushright}
	\begin{textblock*}{0.5\linewidth}(11cm,18.4cm)%{width}[xt,yt](X,Y)
		\justifying Trabalho a ser apresentado à disciplina \disciplina, ministrada pelo professor \docente, no Curso de \curso, da \universidade.
	\end{textblock*} ~\\[1em]
\end{flushright}
\vfill
\vfill
\vfill
\begin{center}
	\cidade{ }- Dezembro de 2017 
\end{center}
\cleardoublepage
%================================================================
% ---------------- FIM DAS PARTES AUTOMÁTICAS -------------------
%================================================================
% Começe a editar seu trabalho a partir dessa parte. Você poderá 
% criar capítuos, seções e os diversos elementos do Latex.
% Use pacotes adicionais para expandir esse documento.
%================================================================
%----------------------------------------------------------------------------%

%----------------------------------------------------------------------------%
% Descrição das ferramentas pesquisadas.
% * Dividimos em seções por assunto e em cada assunto, uma subseção descreve 
%   a ferramenta a ser apresentada.
% * Existem alguns campos como exemplo, mas podemos acrescentar outros se
%   necessário
%----------------------------------------------------------------------------%

%----------------------------------------------------------------------------%
% Início de seção
\section{Assunto}

%----------------------------------------------------------------------------%
% Início de subseção
\subsection{Nome da Ferramenta}

\subsubsection{Descrição}
Breve descrição da ferramenta

Caso você precise representar um código, você poderá usar esse modelo. Pesquise sobre a classe ``minted'' do Latex para suporte a outras linguagens.

\begin{minted}[linenos=true,frame=lines,fontsize=\footnotesize,bgcolor=fundo-cod]{plpgsql}
SELECT numero_conta, nome_agencia 
FROM banco.conta c 
JOIN banco.agencia a 
ON c.cod_agencia=a.cod_agencia;
\end{minted}

\subsubsection{Formato de dados}
Descrição do formato dos dados, quantidade de dados

\subsubsection{Limite de usuários}
Descrever como é tratado o limite de usuários da ferramenta

\subsubsection{API}
Descrição das APIs da ferramenta

\subsubsection{Extensões}
Descrição dos plugins e/ou extensões de interesse

\subsubsection{Licença}
Descrição da licença da ferramenta

\subsubsection{Hospedagem e segurança de dados}
Descrição das implicações da política de segurança da informação e política de privacidade dos dados

\subsubsection{Observações}
Observações, se necessário

\clearpage
% Fim de Subseção
%----------------------------------------------------------------------------%

%----------------------------------------------------------------------------%
% Início de subseção
\subsection{Nome da Ferramenta}

\subsubsection{Descrição}
Breve descrição da ferramenta

\subsubsection{Formato de dados}
Descrição do formato dos dados, quantidade de dados

\subsubsection{Limite de usuários}
Descrever como é tratado o limite de usuários da ferramenta

\subsubsection{API}
Descrição das APIs da ferramenta

\subsubsection{Extensões}
Descrição dos plugins e/ou extensões de interesse

\subsubsection{Licença}
Descrição da licença da ferramenta

\subsubsection{Hospedagem e segurança de dados}
Descrição das implicações da política de segurança da informação e política de privacidade dos dados

\subsubsection{Observações}
Observações, se necessário

\clearpage
% Fim de Subseção
%----------------------------------------------------------------------------%
% Fim de Seção
%----------------------------------------------------------------------------%


%----------------------------------------------------------------------------%
% Início de seção
\section{Assunto seguinte}

%----------------------------------------------------------------------------%
% Início de subseção
\subsection{Nome da Ferramenta}

\subsubsection{Descrição}
Breve descrição da ferramenta

\subsubsection{Formato de dados}
Descrição do formato dos dados, quantidade de dados

\subsubsection{Limite de usuários}
Descrever como é tratado o limite de usuários da ferramenta

\subsubsection{API}
Descrição das APIs da ferramenta

\subsubsection{Extensões}
Descrição dos plugins e/ou extensões de interesse

\subsubsection{Licença}
Descrição da licença da ferramenta

\subsubsection{Hospedagem e segurança de dados}
Descrição das implicações da política de segurança da informação e política de privacidade dos dados

\subsubsection{Observações}
Observações, se necessário

\clearpage
% Fim de Subseção
%----------------------------------------------------------------------------%
% Fim de Seção
%----------------------------------------------------------------------------%
% Informações de licença do documento final (Opcional)
%----------------------------------------------------------------------------%
\vfill{}
\begin{center}
\doclicenseThis
\end{center}
\begin{center}
\textbf{\textcopyleft\docente}{ }- Decom/FACET - UFVJM
\end{center}
%----------------------------------------------------------------------------%

\end{document}